\documentclass[11pt,a4paper]{article}

\usepackage[utf8]{inputenc}
\usepackage[T1]{fontenc}
\usepackage{amsmath,amssymb,amsthm}
\usepackage{mathtools}
\usepackage{hyperref}
\usepackage{cleveref}
\usepackage{booktabs}
\usepackage{graphicx}
\usepackage{xcolor}
\usepackage{listings}
\usepackage[margin=1in]{geometry}
\usepackage{enumitem}
\usepackage{microtype}
\usepackage{natbib}

% Theorem environments
\newtheorem{theorem}{Theorem}[section]
\newtheorem{lemma}[theorem]{Lemma}
\newtheorem{proposition}[theorem]{Proposition}
\newtheorem{corollary}[theorem]{Corollary}
\theoremstyle{definition}
\newtheorem{definition}[theorem]{Definition}
\newtheorem{assumption}[theorem]{Assumption}
\newtheorem{remark}[theorem]{Remark}

% Notation
\newcommand{\R}{\mathbb{R}}
\newcommand{\E}{\mathbb{E}}
\newcommand{\Var}{\mathrm{Var}}
\newcommand{\Cov}{\mathrm{Cov}}
\newcommand{\tr}{\mathrm{tr}}
\newcommand{\proj}[1]{P_{#1}}
\newcommand{\norm}[1]{\left\|#1\right\|}
\newcommand{\abs}[1]{\left|#1\right|}
\newcommand{\inner}[2]{\langle #1, #2 \rangle}
\DeclareMathOperator*{\argmin}{arg\,min}
\DeclareMathOperator*{\argmax}{arg\,max}
\DeclareMathOperator{\im}{im}

% Lean 4 code style
\lstdefinestyle{lean}{
  basicstyle=\ttfamily\footnotesize,
  keywordstyle=\color{blue!70!black}\bfseries,
  commentstyle=\color{gray}\itshape,
  breaklines=true,
  breakatwhitespace=true,
  frame=single,
  framerule=0.4pt,
  rulecolor=\color{gray!50},
  xleftmargin=1em,
  columns=flexible,
  literate={ℝ}{{\ensuremath{\mathbb{R}}}}1
           {ℕ}{{\ensuremath{\mathbb{N}}}}1
           {ℚ}{{\ensuremath{\mathbb{Q}}}}1
           {≤}{{\ensuremath{\leq}}}1
           {≥}{{\ensuremath{\geq}}}1
           {≠}{{\ensuremath{\neq}}}1
           {∀}{{\ensuremath{\forall}}}1
           {∃}{{\ensuremath{\exists}}}1
           {→}{{\ensuremath{\to}}}1
           {←}{{\ensuremath{\leftarrow}}}1
           {↔}{{\ensuremath{\leftrightarrow}}}1
           {∈}{{\ensuremath{\in}}}1
           {∉}{{\ensuremath{\notin}}}1
           {⊆}{{\ensuremath{\subseteq}}}1
           {λ}{{\ensuremath{\lambda}}}1
           {·}{{\ensuremath{\cdot}}}1
           {‖}{{\ensuremath{\|}}}1
           {∨}{{\ensuremath{\lor}}}1
           {∧}{{\ensuremath{\land}}}1
           {⟨}{{\ensuremath{\langle}}}1
           {⟩}{{\ensuremath{\rangle}}}1
           {∑}{{\ensuremath{\sum}}}1
           {Wᗮ}{{$W^\perp$}}2,
  morekeywords={theorem,lemma,def,noncomputable,where,by,have,let,show,
    sorry,Prop,Type,structure,class,instance,section,namespace,open,
    import,variable,example,exact,apply,intro,simp,ring,linarith,
    norm_num,positivity,grind,unfold,rw,fun,if,then,else,split_ifs,
    nlinarith,calc,native_decide},
}

\title{Formalizing Sparse Feature Preservation:\\
13 Lean~4 Theorems and Two Corrections}

\author{
  Diego Quinto\thanks{Correspondence: \texttt{github.com/0xQuinto}. Co-authorship offered to Georgios Konstantopoulos.} \\
  \textit{Independent Researcher}
}

\date{April 2026}

\begin{document}

\maketitle

\begin{abstract}
We formalize 13 theorems in Lean~4 for Sparse Feature Preservation~(SFP), a regularization method that mitigates catastrophic forgetting in continual LLM fine-tuning. The formalization process surfaced two substantive errors in our own informal reasoning: the exponent in an adaptive regularization schedule (corrected from $\exp(-\alpha t/2)$ to $\exp(-\alpha t)$) and the direction of optimal multi-task ordering (corrected from descending to ascending forgetting constant). All 13 theorems compile with zero \texttt{sorry} statements and extend the existing SFP Lean~4 library from 17 to 30 machine-verified results, covering finite-sample estimation bounds, distributional derivation of the core hypothesis, optimizer convergence, multi-task extensions, LoRA confounds, and evaluation axiomatics.
\end{abstract}

%=============================================================================
\section{Introduction}
\label{sec:intro}

Formal verification of mathematical arguments is not bookkeeping---it reveals errors that survive human review. This paper reports two such corrections arising from mechanized proof of informal theorems about Sparse Feature Preservation~(SFP).

SFP~\citep{sfp2025} prevents catastrophic forgetting in LLM fine-tuning by penalizing activation drift in a PCA-identified subspace. Given pre-trained parameters $\theta_0$ and a set of monitored layers, the training objective is
\begin{equation}
\label{eq:sfp-loss}
L(\theta) = L_{\mathrm{new}}(\theta) + \lambda \sum_\ell \norm{\proj{W_\ell}\!\big(a_\ell(\theta) - a_\ell(\theta_0)\big)}^2,
\end{equation}
where $a_\ell(\theta) \in \R^d$ is the mean-pooled activation at layer $\ell$ and $\proj{W_\ell}$ is the orthogonal projection onto a $r$-dimensional subspace identified by PCA. The existing Lean~4 formalization provides 17 theorems establishing that if a ``forgetting constant'' $C > 0$ and preservation loss $\delta$ satisfy hypothesis H1, then $\text{forgetting} \leq C\sqrt{\delta}$.

We contribute 13 additional theorems, bringing the total to 30 machine-verified results. Two of these theorems corrected errors in our prior informal analysis:

\begin{enumerate}[nosep]
\item \textbf{Adaptive $\lambda$ schedule (T6).} We conjectured the optimal regularization schedule was $\lambda(t) = \lambda_0 \exp(-\alpha t / 2)$. The Lean proof of the drift-balance condition forced a correction: the unique schedule achieving zero instantaneous drift is $\lambda(t) = \lambda_0 \exp(-\alpha t)$. The factor of $1/2$ in the exponent is erroneous.

\item \textbf{Multi-task ordering (T9).} We stated that optimal task ordering is descending $C_k$ (highest forgetting constant first). The mechanized exchange argument showed the opposite: ascending $C_k$ is optimal, yielding a 46\% reduction in the forgetting bound in our numerical example.
\end{enumerate}

The 13 theorems span four tiers: bounding $C$ (\S\ref{sec:tier1}), deriving H1 (\S\ref{sec:tier2}), optimizer dynamics (\S\ref{sec:tier3}), and extensions (\S\ref{sec:tier4}). All proofs were generated using Aristotle~\citep{aristotle2025} and are available at:

\smallskip
\noindent\url{https://github.com/0xQuinto/sfp/tree/formal-improvements/lean-proofs}

%=============================================================================
\section{Background}
\label{sec:background}

\paragraph{Notation.} Let $\theta \in \R^p$ denote network parameters, $a_\ell(\theta; x) \in \R^d$ the mean-pooled activation at layer~$\ell$ on input~$x$, and $\Delta a_\ell = a_\ell(\theta) - a_\ell(\theta_0)$ the activation drift. For a subspace $W \subseteq \R^d$ with orthonormal basis $U \in \R^{d \times r}$, the projection is $\proj{W} = UU^\top$.

\paragraph{SFP setup.} At each task boundary, SFP collects activations from a memory buffer of $n$ old-task examples, computes a PCA basis at selected layers, and caches anchor projections. The upstream implementation selects layers at fixed relative depths $\{0.25, 0.5, 0.75\}$ with a default buffer of $n = 128$.

\paragraph{Hypothesis H1.} The core SFP hypothesis posits subspaces $W_\ell$ and a constant $C > 0$ such that $\abs{L_{\mathrm{old}}(\theta) - L_{\mathrm{old}}(\theta_0)} \leq C \sum_\ell \norm{\proj{W_\ell}(\Delta a_\ell)}$. The upstream formalization treats H1 as an axiom; we derive it as a theorem in \S\ref{sec:tier2}.

\paragraph{Formalization workflow.} Each theorem is written as a structured input file (definitions, assumptions, claims), submitted to Aristotle~\citep{aristotle2025}, which generates a complete Lean~4 project with Mathlib dependencies. We verify zero \texttt{sorry} statements.

%=============================================================================
\section{Tier 1: Bounding the Forgetting Constant}
\label{sec:tier1}

\subsection{T1: Forgetting Constant as OLS Coefficient}

The forgetting constant $C_\ell$ has the same structure as Kyle's Lambda in market microstructure~\citep{kyle1985}: both are OLS slopes of the form $\beta = \Cov(Y, X) / \Var(X)$.

\begin{definition}[Forgetting constant]
\label{def:C}
Let $D_\ell = \norm{\proj{W}(\Delta a_\ell)}$ and $F = L_{\mathrm{old}}(\theta) - L_{\mathrm{old}}(\theta_0)$. The forgetting constant is $C_\ell = \Cov(F, D_\ell) / \Var(D_\ell)$.
\end{definition}

\begin{theorem}[OLS identity --- T1]
\label{thm:ols}
Under the linear model $F = C_\ell D_\ell + \varepsilon$ with $\E[\varepsilon D_\ell] = 0$ and $\Var(D_\ell) \neq 0$:
\[
C_\ell = \frac{\E[FD_\ell] - \E[F]\E[D_\ell]}{\E[D_\ell^2] - (\E[D_\ell])^2}.
\]
\end{theorem}

\noindent\textit{Lean~4 (verified, 0 \texttt{sorry}):}
\begin{lstlisting}[style=lean]
theorem ols_coefficient_identity
    (C E_D E_D2 E_F E_FD : ℝ)
    (hVarD : E_D2 - E_D ^ 2 ≠ 0)
    (hFD : E_FD = C * E_D2) -- E[FD] = C · E[D²]
    (hF : E_F = C * E_D) :  -- E[F] = C · E[D]
    C = (E_FD - E_F * E_D) / (E_D2 - E_D ^ 2) := by
  grind
\end{lstlisting}

The \texttt{grind} closer is deliberate: the result is algebraically immediate once stated correctly. The difficulty is \emph{identifying} $C$ as an OLS coefficient, not proving the identity.

\begin{theorem}[Finite-sample concentration --- T1]
\label{thm:concentration}
Under sub-Gaussian assumptions ($D_\ell$ is $\sigma_D$-sub-Gaussian, $F$ is $\sigma_F$-sub-Gaussian, $\Var(D_\ell) \geq v_{\min} > 0$), for $n$ i.i.d.\ samples and $\delta \in (0,1)$:
\[
\Pr\!\left(\abs{\hat{C}_\ell - C_\ell} > \varepsilon\right) \leq \delta, \quad \varepsilon = \frac{\sigma_F \sigma_D}{v_{\min}} \sqrt{\frac{2\log(2/\delta)}{n}}.
\]
The bound is monotone decreasing in $n$ and non-vacuous when $n > 2\sigma_F^2 \sigma_D^2 \log(2/\delta) / (v_{\min}^2 \varepsilon_{\mathrm{target}}^2)$.
\end{theorem}

\begin{theorem}[Layer selection optimality --- T1]
\label{thm:layer-select}
Selecting the $K$ layers with smallest $\hat{C}_\ell$ minimizes $\sum_{\ell \in S} C_\ell \sqrt{\delta_\ell}$ over all $K$-element subsets.
\end{theorem}

\subsection{T2: Davis-Kahan Buffer Sizing}

The PCA subspace $\hat{W}$ estimated from $n$ samples satisfies the Davis-Kahan bound~\citep{daviskahan1970}: $\sin\theta(\hat{W}, W) \leq \norm{\hat{\Sigma} - \Sigma}_{\mathrm{op}} / \Delta_k$, where $\Delta_k = \lambda_k - \lambda_{k+1}$ is the eigenvalue gap. Combined with sub-Gaussian concentration, this gives a minimum buffer size $n \geq \sigma^4 \log(2/\delta) / (c \cdot \Delta_k^2 \cdot t^2)$ for $\sin\theta \leq t$ with probability $\geq 1-\delta$. The preservation loss error from using $\hat{W}$ instead of $W$ is bounded by $\abs{\hat{L}_{\mathrm{pres}} - L_{\mathrm{pres}}} \leq 2\sin\theta \cdot \norm{\Delta a}^2$.

%=============================================================================
\section{Tier 2: Strengthening H1}
\label{sec:tier2}

\subsection{T3: H1 as Theorem Under Exponential Family}

The upstream formalization treats H1 as an axiom. We derive it as a consequence of a distributional assumption.

\begin{assumption}[Exponential family activations]
\label{asm:expfam}
The activation distribution belongs to an exponential family with linear sufficient statistic $T(a_\ell) = \proj{W}(a_\ell)$.
\end{assumption}

\begin{theorem}[Sufficiency of projected activations --- T3]
\label{thm:sufficiency}
Under \Cref{asm:expfam}, $L_{\mathrm{old}}(\theta) = \E_x[g(\proj{W}(a_\ell(\theta; x)))]$ for some function~$g$. The old-task loss depends on activations only through $\proj{W}(a_\ell)$.
\end{theorem}

This follows from the Neyman-Fisher factorization theorem: the likelihood factors through the sufficient statistic, so any loss depending on the likelihood depends only on $T(a_\ell) = \proj{W}(a_\ell)$.

\begin{theorem}[H1 as consequence --- T3]
\label{thm:h1}
Under \Cref{asm:expfam} with $g$ being $C$-Lipschitz:
$\abs{L_{\mathrm{old}}(\theta) - L_{\mathrm{old}}(\theta_0)} \leq C \cdot \norm{\proj{W}(\Delta a_\ell)}.$
\end{theorem}

\noindent\textit{Lean~4 (verified, 0 \texttt{sorry}):}
\begin{lstlisting}[style=lean]
theorem H1_from_lipschitz
    (W : Submodule ℝ E) [W.HasOrthogonalProjection]
    (g : E → ℝ) (C : ℝ≥0) (hLip : LipschitzWith C g)
    (L_old : E → ℝ)
    (hfactor : ∀ a : E, L_old a = g (projOnto W a))
    (a a₀ : E) :
    |L_old a - L_old a₀| ≤ C * ‖projOnto W (a - a₀)‖ := by
  -- proof in appendix
\end{lstlisting}

\begin{corollary}[Orthogonal complement irrelevance --- T3]
Under the same assumptions, $\partial L_{\mathrm{old}} / \partial(\proj{W^\perp}(a_\ell)) = 0$. Fine-tuning can freely modify the orthogonal complement. The Lean formalization verifies: for any $v \in W^\perp$, $L_{\mathrm{old}}(a + v) = L_{\mathrm{old}}(a)$.
\end{corollary}

\begin{remark}[Diagnostic]
The assumption is testable via $R^2 = 1 - \Var(L_{\mathrm{old}} \mid \proj{W}(a_\ell)) / \Var(L_{\mathrm{old}})$. The empirical $R^2 = 0.933$ for Qwen2.5-1.5B~\citep{sfp2025} supports it.
\end{remark}

\subsection{T4: Ablation Importance Equals Causal Importance}

Under an additivity assumption ($f(\alpha_1, \ldots, \alpha_d) = g_0 + \sum_i g_i(\alpha_i)$), the ablation importance of each PCA direction equals its Shapley value~\citep{shapley1953}. The ablation ranking is therefore the unique fair ranking. Without additivity, we provide a counterexample where the ablation-top-$r$ set is suboptimal by a factor of $\Omega(d)$. The Lean formalization includes both the equivalence proof and the explicit counterexample construction.

\subsection{T5: PCA vs.\ Gradient-SVD Optimality}

The optimal preservation subspace minimizing expected squared forgetting $\E[(\nabla L^\top \proj{W^\perp} \Delta a)^2]$ is the top-$r$ eigenspace of $M = \Sigma^{1/2} G \Sigma^{1/2}$, where $\Sigma$ is the drift covariance and $G$ the gradient outer product. This follows from the Ky Fan inequality. PCA ($G \propto I$) and gradient-SVD ($\Sigma \propto I$) emerge as special cases. A computable decision score determines which is preferable.

%=============================================================================
\section{Tier 3: Optimizer Dynamics}
\label{sec:tier3}

\subsection{T6: The Adaptive $\lambda$ Correction}
\label{sec:t6}

This is the paper's primary exhibit of formalization catching an error.

\paragraph{Setup.} Suppose the gradient signal in the preserved subspace decays exponentially: $s(t) = \norm{\proj{W}\nabla L_{\mathrm{new}}(\theta_t)}^2 = s_0 e^{-\alpha t}$. The drift evolves as $\delta(t+1) \approx (1 - 2\eta\lambda(t))\delta(t) + \eta \cdot s(t)$. We sought a schedule $\lambda(t)$ that holds $\delta(t)$ at a constant target $\delta_{\mathrm{target}}$.

\paragraph{The conjecture.} Setting $d\delta/dt = 0$ with $\delta = \delta_{\mathrm{target}}$ requires $2\lambda(t) \cdot \delta_{\mathrm{target}} = \eta \cdot s(t)$, giving $\lambda(t) = \eta s_0 e^{-\alpha t} / (2\delta_{\mathrm{target}}) = \lambda_0 e^{-\alpha t}$. We initially wrote $\lambda_0 \exp(-\alpha t / 2)$.

\paragraph{How Lean caught it.} The proof obligation for zero instantaneous drift reduces to showing that $\lambda_0 \exp(-\alpha t / 2) \cdot \delta_{\mathrm{target}}$ cancels $\eta \cdot s_0 \exp(-\alpha t)$. This requires $\exp(-\alpha t / 2) = \exp(-\alpha t)$, which is false. The Lean file makes the correction explicit:

\begin{lstlisting}[style=lean]
-- The original (incorrect) conjecture
def adaptiveLambdaSchedule (lam₀ a : ℝ) (t : ℝ) : ℝ :=
  lam₀ * exp (-a * t / 2)

-- The corrected schedule
def correctAdaptiveSchedule (lam₀ a : ℝ) (t : ℝ) : ℝ :=
  lam₀ * exp (-a * t)
\end{lstlisting}

The theorem that closes the correction:
\begin{lstlisting}[style=lean]
theorem exponential_schedule_constant_drift
    (eta s₀ a d_target : ℝ)
    (_heta : 0 < eta) (_hs₀ : 0 < s₀) (hd : 0 < d_target) :
    let lam₀ := eta * s₀ / (2 * d_target)
    ∀ t : ℝ, instantaneousDriftRate eta (signalDecay s₀ a t)
      (correctAdaptiveSchedule lam₀ a t) d_target = 0 := by
  -- unfold + ring; see appendix
\end{lstlisting}

The corrected schedule uses $\exp(-\alpha t)$, not $\exp(-\alpha t/2)$. The original error was a factor-of-two mistake in tracing through the drift balance equation---exactly the kind of error that persists through informal review but cannot survive a proof checker.

\subsection{T7: SGD Convergence Under SFP}

Under standard assumptions ($L_{\mathrm{new}}$ is $L$-smooth and $\mu$-strongly convex, step size $\eta \leq 1/(L + \lambda L_p)$), after $K$ gradient steps:
\begin{align}
L_{\mathrm{new}}(\theta_K) - L_{\mathrm{new}}(\theta^*) &\leq \left(1 - \frac{\mu}{L + \lambda L_p}\right)^K\!\! (L_{\mathrm{new}}(\theta_0) - L_{\mathrm{new}}(\theta^*)) + \frac{\lambda \delta_{\max}}{\mu}, \label{eq:convergence} \\
\delta_K &\leq (1 - \eta\mu)^K \delta_0 + \frac{\eta\lambda}{\mu}\norm{\nabla L_{\mathrm{pres}}}_{\max}^2. \label{eq:drift}
\end{align}
The effective condition number is $\kappa_{\mathrm{eff}} = (L + \lambda L_p)/\mu$. Eliminating $\lambda$ yields a Pareto frontier: $\text{forgetting} \leq C \cdot (L_{\mathrm{new}} - L_{\mathrm{new}}(\theta^*))^{1/4} \cdot (\eta \norm{\nabla L_{\mathrm{pres}}}_{\max}^2 / \mu^2)^{1/4}$.

\subsection{T8: Bound Tightness}

The exponent $1/2$ in $\text{forgetting} \leq C\sqrt{\delta}$ is tight: the construction $L_{\mathrm{old}}(a) = \norm{\proj{W}(a - a_0)}$ achieves $\text{forgetting} = \sqrt{\delta}$ exactly. However, when $\nabla L_{\mathrm{old}}(\theta_0) = 0$ (the typical LLM regime where $\theta_0$ is near-optimal on the old task), forgetting scales as $O(\delta)$---the bound is loose by $\sqrt{\delta}$. This dichotomy explains why SFP works better in practice than the worst-case bound suggests.

%=============================================================================
\section{Tier 4: Extensions}
\label{sec:tier4}

\subsection{T9: The Multi-Task Ordering Correction}
\label{sec:t9}

This is the paper's second exhibit of formalization catching an error.

\paragraph{Setup.} Consider sequential fine-tuning on tasks $T_1, \ldots, T_K$. The cumulative forgetting bound telescopes as $F_{\mathrm{total}} \leq \sum_{k=1}^{K-1} C_{\pi(k)} \cdot w_k$, where $w_k = (K-k)\sqrt{\delta_{\max}}$ is a decreasing weight (later tasks incur less subsequent drift).

\paragraph{The correction.} To minimize $\sum_k C_{\pi(k)} w_k$ with $w_k$ decreasing, the rearrangement inequality says: pair the largest weight with the \emph{smallest} $C$. The optimal order is \emph{ascending} $C_k$---fine-tune the least forgettable task first. We initially stated the opposite.

The Lean formalization verifies the exchange argument:
\begin{lstlisting}[style=lean]
theorem exchange_reduces_cost
    (n : ℕ) (C : Fin n → ℝ) (w : Fin n → ℝ)
    (hw : ∀ i j : Fin n, i < j → w j < w i) -- w decreasing
    (i : Fin n) (hi : i.val + 1 < n)
    (hCi : C i > C ⟨i.val + 1, hi⟩) :
    -- swapping adjacent out-of-order pair reduces cost
    ∑ k, C' k * w k < ∑ k, C k * w k := by
\end{lstlisting}

\paragraph{Numerical verification.} With $K = 5$, $C \in \{0.5, 0.8, 1.2, 0.3, 0.7\}$, $\sqrt{\delta_{\max}} = 0.1$: ascending~$C$ gives bound~$0.49$ vs.\ descending~$0.91$ (46\% reduction). All three numerical claims are machine-checked by \texttt{norm\_num}.

\subsection{T10: Per-Layer Optimal Weighting}

Minimizing $\sum_\ell C_\ell \sqrt{s_\ell / (2\lambda_\ell)}$ subject to $\sum_\ell \lambda_\ell = \Lambda$ yields, via Lagrange multipliers, the optimal weights $\lambda_\ell^* = \Lambda \cdot (C_\ell \sqrt{s_\ell})^{2/3} / \sum_j (C_j \sqrt{s_j})^{2/3}$, where $s_\ell = \norm{\proj{W_\ell}\nabla_{a_\ell} L_{\mathrm{new}}}^2$. When gradient signal is uniform across layers: $\lambda_\ell^* \propto C_\ell^{2/3}$. The Lean proof establishes $F(\lambda^*) \leq F(\lambda_{\mathrm{uniform}})$ via a tangent-line convexity argument.

\subsection{T11: Projection Information Loss}

Under Gaussian assumptions ($a \sim \mathcal{N}(\mu, \Sigma)$, $Y = w^\top a + \varepsilon$), the information lost by projecting to the top-$r$ PCA subspace is $I(Y; a) - I(Y; \proj{W}(a)) = \frac{1}{2}\log(1 + \norm{\proj{W^\perp} w}^2 \lambda_{r+1} / \sigma_\varepsilon^2)$. When task-relevant directions align with high-variance PCA directions, the explained variance ratio $\rho(r)$ approximates the fraction of task information retained.

\subsection{T12: LoRA Confound}

Under LoRA~\citep{hu2022lora}, activation drift $\Delta a = BAx$ lies in $\im(B)$ with $\dim(\im(B)) \leq r_{\mathrm{lora}}$. When $r_{\mathrm{lora}} \leq r_{\mathrm{sfp}}$ and $\im(B) \subseteq W_{\mathrm{SFP}}$, H1 holds trivially---all drift is captured by the preservation subspace by construction, not by learned feature structure. The high $R^2$ from drift-forgetting regression is an artifact of low-rank confinement. The confound is negligible when $r_{\mathrm{lora}} \geq 2r_{\mathrm{sfp}}$. This theorem is a methodological warning: full fine-tuning experiments are the gold standard for validating~H1.

\subsection{T13: Harmonic Mean Scorer}

The leaderboard scorer $H(R, P) = 2RP/(R+P)$ is the unique power mean satisfying zero-forcing ($H = 0$ iff $R = 0$ or $P = 0$). It equals the Nash bargaining solution~\citep{nash1950} between retention-maximizing and plasticity-maximizing agents. The failure of the previously used weighted linear score is machine-checked:

\begin{lstlisting}[style=lean]
theorem weightedLinear_violates_P1 :
    weightedLinear 0 1 = 0.4 := by norm_num
\end{lstlisting}

A model that forgets everything scores $0.4$ under the linear scorer but $0$ under the harmonic mean.

%=============================================================================
\section{Discussion}
\label{sec:discussion}

\paragraph{What class of error does formalization catch?} Both corrections involve a sign or factor error in an exponent (T6) or inequality direction (T9) that passed informal review. The T6 error requires working through the full drift balance equation to spot; the T9 error requires carefully tracking which direction the rearrangement inequality runs. Lean~4 forces complete proof terms that cannot silently absorb such errors.

\paragraph{What formalization does not do.} The Lean proofs verify that abstract mathematical claims follow from stated assumptions. They do not verify that the assumptions hold for real LLMs, that numerical constants ($\sigma_F$, $\sigma_D$, $v_{\min}$) match practice, or that the Python implementation correctly reflects the theorems. This gap is the main limitation.

\paragraph{Framing these corrections.} Neither correction invalidates the SFP algorithm. T6 tightens the adaptive schedule (which is not yet implemented upstream). T9 corrects the ordering heuristic for multi-task settings (also not yet implemented). The corrections are valuable precisely because they were found \emph{before} implementation, not after a production failure.

\paragraph{Future work.} (a)~Close the T7 Pareto frontier to explicit constants for specific architectures. (b)~Formalize the adaptive $\alpha$ estimation error from T6 more tightly. (c)~Investigate whether T12's LoRA confound can be turned into a positive result (conditions under which LoRA~+~SFP is provably redundancy-free).

%=============================================================================
\section{Conclusion}
\label{sec:conclusion}

We contribute 13 machine-checked Lean~4 theorems strengthening the mathematical foundations of SFP, spanning layer selection, optimizer dynamics, multi-task extensions, and evaluation design. The formalization process caught two substantive errors in our prior informal reasoning: the exponent in the adaptive $\lambda$ schedule~(T6) and the direction of the optimal task ordering~(T9). All proofs are zero-\texttt{sorry} and available at the project repository.

\bibliographystyle{plainnat}
\bibliography{references}

%=============================================================================
\appendix

\section{Theorem Correspondence Table}
\label{app:table}

\begin{table}[h]
\centering
\footnotesize
\begin{tabular}{@{}cllll@{}}
\toprule
\textbf{ID} & \textbf{Name} & \textbf{Lean File} & \textbf{Key Lemma} & \textbf{Sorries} \\
\midrule
T1 & Forgetting constant & \texttt{ForgettingConstantBound.lean} & \texttt{ols\_coefficient\_identity} & 0 \\
T2 & Davis-Kahan buffer & \texttt{DavisKahan.lean} & \texttt{buffer\_size\_sufficient} & 0 \\
T3 & Conditional H1 & \texttt{T3ConditionalH1.lean} & \texttt{H1\_from\_lipschitz} & 0 \\
T4 & Ablation = Shapley & \texttt{AblationCausalImportance.lean} & \texttt{ablation\_eq\_shapley} & 0 \\
T5 & PCA vs gradient-SVD & \texttt{T5PcaVsGradientSvd.lean} & \texttt{combined\_optimality} & 0 \\
T6 & Adaptive $\lambda$ & \texttt{T6AdaptiveLambda\ldots.lean} & \texttt{exponential\_schedule\ldots} & 0 \\
T7 & SGD convergence & \texttt{SGDConvergence.lean} & \texttt{linear\_convergence\ldots} & 0 \\
T8 & Bound tightness & \texttt{BoundTightness.lean} & \texttt{forgetting\_eq\_sqrt\ldots} & 0 \\
T9 & Multi-task ordering & \texttt{T9MultiTask\ldots.lean} & \texttt{exchange\_reduces\_cost} & 0 \\
T10 & Optimal weighting & \texttt{T10MultiLayer\ldots.lean} & \texttt{optWeight\_minimizes} & 0 \\
T11 & Information loss & \texttt{T11Projection\ldots.lean} & \texttt{infoLoss\_le} & 0 \\
T12 & LoRA confound & \texttt{T12LoRAConfound.lean} & \texttt{out\_of\_SFP\_sufficient\ldots} & 0 \\
T13 & Harmonic mean & \texttt{HarmonicMeanScorer.lean} & \texttt{harmonicMean\_zero\_iff} & 0 \\
\bottomrule
\end{tabular}
\caption{Correspondence between informal theorems and Lean~4 identifiers. All 13 compile with zero \texttt{sorry}.}
\label{tab:correspondence}
\end{table}

\section{T3 Proof Sketch: Deriving H1}
\label{app:t3}

Under \Cref{asm:expfam}, the activation distribution at layer $\ell$ has density
\[
p(a_\ell \mid x) = h(a_\ell) \exp\!\big(\eta(x)^\top T(a_\ell) - A(\eta(x))\big)
\]
with $T(a_\ell) = \proj{W}(a_\ell)$.

\textit{Step 1 (Sufficiency).} By the Neyman-Fisher factorization theorem, the likelihood ratio $p(a_\ell \mid x) / p(a_\ell \mid x')$ depends on $a_\ell$ only through $T(a_\ell) = \proj{W}(a_\ell)$. Any loss function $L_{\mathrm{old}}(\theta)$ that depends on the distribution of $a_\ell$ therefore factors as $L_{\mathrm{old}}(\theta) = g(\proj{W}(a_\ell(\theta)))$ for some $g$.

\textit{Step 2 (Lipschitz bound).} If $g$ is $C$-Lipschitz ($\abs{g(u) - g(v)} \leq C\norm{u-v}$ for all $u, v$), then $\abs{L_{\mathrm{old}}(\theta) - L_{\mathrm{old}}(\theta_0)} = \abs{g(\proj{W}(a_\ell(\theta))) - g(\proj{W}(a_\ell(\theta_0)))} \leq C \norm{\proj{W}(\Delta a_\ell)}$, which is H1.

\textit{Step 3 (Orthogonal irrelevance).} Since $g$ depends only on $\proj{W}(a_\ell)$, adding any $v \in W^\perp$ to $a_\ell$ does not change the loss: $\proj{W}(a_\ell + v) = \proj{W}(a_\ell)$.

\textit{Step 4 (Diagnostic).} If $R^2 \coloneqq 1 - \Var(L_{\mathrm{old}} \mid \proj{W}(a_\ell)) / \Var(L_{\mathrm{old}}) = 1$, then $L_{\mathrm{old}}$ is a deterministic function of $\proj{W}(a_\ell)$, which implies the factorization in Step~1. The empirical $R^2 = 0.933$ indicates the assumption is approximately satisfied.

\end{document}
